\documentclass{article}
\usepackage{amsmath,amssymb,xurl}
\usepackage[hidelinks]{hyperref}
% Shared commands for notes in docs/paper-gaps/.

\DeclareUnicodeCharacter{2080}{\ifmmode{}_{0}\else\textsubscript{0}\fi}
\DeclareUnicodeCharacter{2081}{\ifmmode{}_{1}\else\textsubscript{1}\fi}
\DeclareUnicodeCharacter{2082}{\ifmmode{}_{2}\else\textsubscript{2}\fi}
\DeclareUnicodeCharacter{2099}{\ifmmode{}_{n}\else\textsubscript{n}\fi}

\newcommand{\Lean}{\textsc{Lean}}
\ExplSyntaxOn
\NewDocumentCommand{\leanid}{m}{
  \ifmmode
    \textsf{\detokenize{#1}}
  \else
    \group_begin:
    \urlstyle{sf}
    \tl_set:Nn \l_tmpa_tl {#1}
    \tl_replace_all:Nnn \l_tmpa_tl {\_} {_}
    \exp_args:NV \path \l_tmpa_tl
    \group_end:
  \fi
}
\ExplSyntaxOff
\providecommand{\tr}{\operatorname{tr}}
\newcommand{\tnleanIssueBase}{https://github.com/LionSR/TNLean/issues/}
\newcommand{\tnleanPrBase}{https://github.com/LionSR/TNLean/pull/}
\newcommand{\tnleanDocBase}{https://sirui-lu.com/TNLean/docs/}
\newcommand{\ghissue}[1]{\href{\tnleanIssueBase#1}{GitHub issue \##1}}
\newcommand{\ghpr}[1]{\href{\tnleanPrBase#1}{PR \##1}}
\newcommand{\doclink}[1]{\href{\tnleanDocBase#1.html}{\path{#1}}}

% Dirac notation and the identity operator, as in blueprint/src/macros/common.tex.
% \providecommand keeps note-local definitions authoritative.
\usepackage{dsfont}
\providecommand{\ket}[1]{|#1\rangle}
\providecommand{\bra}[1]{\langle #1|}
\providecommand{\braket}[2]{\langle #1 | #2 \rangle}
\providecommand{\ketbra}[2]{|#1\rangle\!\langle #2|}
\providecommand{\Id}{\mathds{1}}
\providecommand{\one}{\mathds{1}}
\providecommand{\C}{\mathbb{C}}
\providecommand{\MN}[1]{M_{#1}(\C)}
\providecommand{\MD}{M_D(\C)}

% Verdict marker: \gapnote{<kind>}{<status>}. Kind is the policy
% classification (clarification, local-correction, scope-restriction,
% unfaithful, false-source, open-gap); status is open, wip (elimination
% underway), resolved, or historical. Machine-read by the site index;
% prints a verdict line.
\newcommand{\gapnote}[2]{\par\noindent\textbf{Verdict:} #1 (#2).\par}

\title{Noncentral Tree Gauges with a Fixed Untwisted Component}
\author{}
\date{}
\begin{document}
\maketitle
\gapnote{scope-restriction}{open}

Schuch, Cirac, and P\'erez-Garc\'ia,
\href{https://arxiv.org/abs/1001.3807v3}{arXiv:1001.3807v3},
Lemma 6.14 and the exterior-block argument of Theorem 6.9, source
\path{Papers/1001.3807/paper_v3.tex}, lines 1935--1990 and 2199--2214,
motivate removing complementary tree operators while keeping the boundary
coordinates fixed. Centrality cannot be discarded from the local Gram
argument: for noncentral insertions a different contraction argument is
required.

The theorem
\leanid{TNLean.PEPS.exists_regularRegionTreeGauge_fixed_on_untwisted_component}
constructs a spanning tree and a root-normalized gauge from a connected
region and a supplied untwisted subgraph. It first takes a spanning forest
of the untwisted subgraph and then extends the forest to a tree of the
region. The resulting gauge fixes every vertex in the root's untwisted
component and removes every internal tree operator. No commutativity,
tree, forest, gauge, factorization or density equality is supplied as a
hypothesis.

The theorem
\leanid{TNLean.PEPS.exists_regularComplementGauge_fixed_on_incident}
then assumes that every complementary boundary endpoint is in the root's
untwisted component, together with identity on every original incident
bond. It extends the derived complementary gauge by the identity on the
observation region, proves preservation of identity on every incident
bond and removal of all complementary tree operators, and derives equality
of the actual closed PEPS vectors from local regular invariance.

These establish noncentral algebraic steps. To apply them to an actual open
flux pair, the physical exterior collar must still be shown to project to
one untwisted component containing every complementary boundary endpoint.
The induced gauge must then be composed with the actual contraction
identities. These theorems do not assume or
assert those geometric and density conclusions, and it does not identify
arbitrary endpoint-free annular observations with the vacuum. The remaining
source-level conclusions stay under \ghissue{8703}. The full-lattice
Hamiltonian work of \ghissue{8947} is unaffected.
\end{document}
